\chapter{Capitolo 17 - Sintesi, agenda 2026-2030, protocollo di adozione scientifica}

\section{1. Problema}

Un testo fondativo non e completo quando termina con un elenco di risultati. E completo quando chiarisce:

\begin{enumerate}
\def\labelenumi{\arabic{enumi}.}
\tightlist
\item
  che cosa e stato realmente dimostrato;
\item
  che cosa resta in revisione;
\item
  come una comunità scientifica può adottare, criticare e migliorare il paradigma senza dipendere da un singolo autore o da un singolo laboratorio.
\end{enumerate}

Per questo il capitolo conclusivo affronta un problema diverso dai capitoli precedenti:

\textbf{come trasformare OCT da architettura teorica e programma interno in infrastruttura scientifica condivisibile, verificabile e cumulativa tra il 2026 e il 2030?}

\subsection{1.1 Cosa deve fare una chiusura fondativa}

Una chiusura seria deve evitare due errori simmetrici:

\begin{enumerate}
\def\labelenumi{\arabic{enumi}.}
\tightlist
\item
  trionfalismo prematuro (presentare come consolidato ciò che e ancora \texttt{in\_proof/revise});
\item
  scetticismo sterile (negare i progressi già strutturati e validati).
\end{enumerate}

La via scelta e una sintesi a tre livelli:

\begin{enumerate}
\def\labelenumi{\arabic{enumi}.}
\tightlist
\item
  livello teorico-formale;
\item
  livello metodologico-validativo;
\item
  livello istituzionale-operativo.
\end{enumerate}

\subsection{1.2 Stato di maturità all'uscita del libro}

Alla data corrente, OCT e in stato di maturità ``fondativo operativo'':

\begin{enumerate}
\def\labelenumi{\arabic{enumi}.}
\tightlist
\item
  lessico, assiomi e mappe classico -\textgreater{} ordinativo sono esplicitati;
\item
  protocollo di validazione e governance sono formalizzati;
\item
  due case study (fisico e sistemico) mostrano applicabilità strutturale;
\item
  persiste una quota di claims in \texttt{in\_review/in\_proof} che richiede cicli ulteriori.
\end{enumerate}

Questo stato e sufficiente per pubblicazione scientifica aperta, non ancora per chiusura definitiva della teoria.

\begin{center}\rule{0.5\linewidth}{0.5pt}\end{center}

\section{2. Tesi del capitolo}

La tesi conclusiva e articolata in nove enunciati.

\begin{enumerate}
\def\labelenumi{\arabic{enumi}.}
\tightlist
\item
  OCT e una estensione necessaria della lettura categoriale quando la sola preservazione strutturale non basta a discriminare validità ordinativa.
\item
  Il cuore innovativo di OCT e la grammatica di coerenza contestuale (\texttt{Coh}, \texttt{Phi}, \texttt{Delta}) e non una semplice rinominazione di categoria classica.
\item
  L'unione di formalismo, benchmark e governance costituisce il minimo per una teoria cumulativa.
\item
  I case study mostrano che la forma ordinativa e trasferibile tra domini con diversa natura epistemica.
\item
  Il programma resta scientifico solo se mantiene pubblici limiti, fallimenti e stati di revisione.
\item
  La fase 2026-2030 deve essere guidata da milestone verificabili e non da dichiarazioni di principio.
\item
  La disseminazione multi-piattaforma ha valore solo con sincronizzazione rigorosa di versione e stato.
\item
  L'adozione internazionale richiede una didattica a bassa soglia e una validazione ad alta soglia.
\item
  Il successo del programma dipende dalla capacità di attrarre confutazioni qualificate, non solo adesioni.
\end{enumerate}

\begin{center}\rule{0.5\linewidth}{0.5pt}\end{center}

\section{3. Bilancio sintetico dell'opera}

\subsection{3.1 Contributi principali consolidati}

L'opera ha prodotto sei contributi strutturali.

\begin{enumerate}
\def\labelenumi{\arabic{enumi}.}
\tightlist
\item
  \textbf{Fondazione ontologica ordinativa}: passaggio da oggetti statici a singolarità relazionali e funzioni emergenti.
\item
  \textbf{Estensione categoriale}: reinterpretazione di nozioni classiche con vincoli di non-degenerazione e coerenza contestuale.
\item
  \textbf{Metateoria della validazione}: protocollo PV-8, criteri di evidenza minima, decisione tracciata.
\item
  \textbf{Governance epistemica}: doppio stato claim, quality gates, audit trail, gestione eccezioni.
\item
  \textbf{Case study fisico}: controfase strutturale e quantizzazione topologica in sistema replicabile.
\item
  \textbf{Case study sistemico}: modello reaction-diffusion multi-scala con biforcazione tripla e distinzione \texttt{senza\ vista/con\ vista}.
\end{enumerate}

\subsection{3.2 Risultati in consolidamento}

Restano in consolidamento:

\begin{enumerate}
\def\labelenumi{\arabic{enumi}.}
\tightlist
\item
  promozione theorem-level di claims ad alta portata metateorica;
\item
  benchmarking indipendente esteso per dominio sociale;
\item
  stabilizzazione di metriche ausiliarie cross-dominio.
\end{enumerate}

Questa zona non e debolezza narrativa: e il luogo normale della ricerca viva.

\subsection{3.3 Cosa non e OCT}

Per chiarezza pubblica, OCT non e:

\begin{enumerate}
\def\labelenumi{\arabic{enumi}.}
\tightlist
\item
  un sostituto totale della teoria delle categorie classica;
\item
  un sistema predittivo onnipotente;
\item
  una cornice ideologica non falsificabile.
\end{enumerate}

OCT e una grammatica di composizione reale con regole di verifica esplicite.

\subsection{3.4 Dialogo con letteratura adiacente (perimetro non autoreferenziale)}

Per robustezza scientifica, OCT deve essere sempre letta in dialogo con linee di ricerca gia attive.

Matrice minima di confronto:

\begin{enumerate}
\def\labelenumi{\arabic{enumi}.}
\item
  \textbf{Applied Category Theory (ACT)}

  \begin{itemize}
  \tightlist
  \item
    contributo esterno: modellazione composizionale di sistemi e database;
  \item
    contributo OCT: gate di ammissibilita ontologica (\texttt{Coh/Phi/Delta}) e stato claim tracciato.
  \end{itemize}
\item
  \textbf{Process theories / Categorical Quantum Mechanics}

  \begin{itemize}
  \tightlist
  \item
    contributo esterno: formalizzazione dei processi e della composizione fisica;
  \item
    contributo OCT: filtro di non-degenerazione contestuale sotto perturbazione.
  \end{itemize}
\item
  \textbf{Open systems e dynamical systems compositionali}

  \begin{itemize}
  \tightlist
  \item
    contributo esterno: scambio con ambiente, dinamiche non chiuse;
  \item
    contributo OCT: integrazione nativa con protocollo validativo e governance di evidenza.
  \end{itemize}
\item
  \textbf{Resource theories}

  \begin{itemize}
  \tightlist
  \item
    contributo esterno: convertibilita, monotoni, costo trasformativo;
  \item
    contributo OCT: distinzione esplicita tra struttura composibile e funzione emergente non sterile.
  \end{itemize}
\item
  \textbf{Resilience / complexity science}

  \begin{itemize}
  \tightlist
  \item
    contributo esterno: regimi di stabilita, transizioni, attrattori;
  \item
    contributo OCT: innesto categoriale + audit trail dei passaggi decisionali.
  \end{itemize}
\end{enumerate}

Regola pratica:
ogni benchmark OCT di livello \texttt{L2+} deve dichiarare almeno una baseline esterna della matrice sopra, per evitare chiusura autoreferenziale.

\begin{center}\rule{0.5\linewidth}{0.5pt}\end{center}

\section{4. Agenda 2026-2030}

\subsection{4.1 Obiettivo generale}

Portare OCT da baseline fondativa \texttt{v1.x} a standard scientifico interoperabile \texttt{v2.0} entro il 2030.

\subsection{4.2 Roadmap per fasi}

\textbf{Fase A (2026-2027): Consolidamento metodologico}

\begin{enumerate}
\def\labelenumi{\arabic{enumi}.}
\tightlist
\item
  completamento cicli benchmark indipendenti su claim \texttt{in\_proof} prioritari;
\item
  hardening del protocollo PV-8 con template dominio-specifici;
\item
  rilascio di dataset/manifest con pacchetti riproducibili \texttt{R2/R3}.
\end{enumerate}

\textbf{Fase B (2027-2028): Espansione disciplinare}

\begin{enumerate}
\def\labelenumi{\arabic{enumi}.}
\tightlist
\item
  ingresso di almeno tre domini aggiuntivi (es. biologia dei pattern, sistemi economici, reti cognitive);
\item
  definizione di ontologie ponte per interoperabilità notazionale;
\item
  pubblicazioni comparative con baseline non ordinative.
\end{enumerate}

\textbf{Fase C (2028-2029): Validazione inter-laboratorio}

\begin{enumerate}
\def\labelenumi{\arabic{enumi}.}
\tightlist
\item
  repliche indipendenti gestite da team esterni;
\item
  audit incrociato di decision matrix e claim ledger;
\item
  benchmark challenge pubblica con protocolli preregistrati.
\end{enumerate}

\textbf{Fase D (2029-2030): Standardizzazione e formazione}

\begin{enumerate}
\def\labelenumi{\arabic{enumi}.}
\tightlist
\item
  emissione OCT \texttt{v2.0} con specifica stabile;
\item
  handbook didattico modulare IT/EN;
\item
  rete di adozione scientifica con criteri di qualità certificabili.
\end{enumerate}

\subsection{4.3 KPI strategici 2026-2030}

Indicatori minimi proposti:

\begin{enumerate}
\def\labelenumi{\arabic{enumi}.}
\tightlist
\item
  percentuale claim con ciclo L1-L2-L3 completo;
\item
  numero repliche indipendenti riuscite;
\item
  tempo medio di chiusura \texttt{revise} con decisione tracciata;
\item
  copertura dominio dei benchmark pubblici;
\item
  tasso di coerenza multi-piattaforma release.
\end{enumerate}

Soglie KPI devono essere pubbliche e versionate annualmente.

\begin{center}\rule{0.5\linewidth}{0.5pt}\end{center}

\section{5. Protocollo di adozione scientifica}

\subsection{5.1 Livelli di adozione}

Proponiamo quattro livelli progressivi.

\begin{enumerate}
\def\labelenumi{\arabic{enumi}.}
\tightlist
\item
  \texttt{L0\ -\ Lettura}: uso concettuale senza impegni validativi.
\item
  \texttt{L1\ -\ Implementazione\ locale}: applicazione su casi interni con report standard.
\item
  \texttt{L2\ -\ Replica\ aperta}: pubblicazione pacchetto riproducibile e confronto baseline.
\item
  \texttt{L3\ -\ Nodo\ di\ rete}: laboratorio che contribuisce a benchmark, audit e governance comune.
\end{enumerate}

\subsection{\texorpdfstring{5.2 Requisiti minimi per ingresso \texttt{L2}}{5.2 Requisiti minimi per ingresso L2}}

\begin{enumerate}
\def\labelenumi{\arabic{enumi}.}
\tightlist
\item
  adozione del protocollo PV-8;
\item
  claim ledger esplicito \texttt{S0-S3};
\item
  quality gates \texttt{G0-G6} pubblicati;
\item
  rilascio artefatti con manifest e hash;
\item
  dichiarazione limiti/non-claim nel report.
\item
  confronto esplicito con almeno una baseline esterna (ACT/process/open systems/resource/resilience), con criterio di superiorita o complementarita dichiarato.
\end{enumerate}

\subsection{5.3 Regole etiche ed epistemiche}

Per evitare abuso del framework:

\begin{enumerate}
\def\labelenumi{\arabic{enumi}.}
\tightlist
\item
  divieto di presentare claim \texttt{in\_review} come risultati consolidati;
\item
  obbligo di pubblicare errori significativi insieme ai successi;
\item
  tracciabilità delle eccezioni governance;
\item
  separazione tra analisi scientifica e uso narrativo/politico.
\end{enumerate}

\subsection{5.4 Percorso didattico consigliato}

Per adozione internazionale graduale:

\begin{enumerate}
\def\labelenumi{\arabic{enumi}.}
\tightlist
\item
  modulo ``Start here'' per lettori non specialisti;
\item
  modulo formalismo (assiomi, notazione, teoremi);
\item
  modulo laboratorio (benchmark, metriche, decisione);
\item
  modulo governance (release, audit, versioning).
\end{enumerate}

L'obiettivo e fruizione semplice in ingresso e rigore forte in uscita.

\begin{center}\rule{0.5\linewidth}{0.5pt}\end{center}

\section{6. Proposizioni e teoremi conclusivi}

\subsection{Proposizione 17.1 (Sufficienza fondativa v1.x)}

Enunciato:
L'insieme dei capitoli 1-16 e sufficiente per qualificare OCT come programma fondativo scientificamente pubblicabile.

Intuizione:
sono presenti teoria, metodo, governance e casi applicativi con stati dichiarati.

Stato: \texttt{validated}.

\subsection{Proposizione 17.2 (Necessità della fase 2026-2030)}

Enunciato:
Senza roadmap benchmark inter-laboratorio, OCT resta cornice promettente ma non standard disciplinare.

Intuizione:
la scientificità cumulativa richiede replica esterna e confronto competitivo.

Stato: \texttt{validated}.

\subsection{\texorpdfstring{Teorema 17.3 (Condizione di transizione a standard \texttt{v2.0})}{Teorema 17.3 (Condizione di transizione a standard v2.0)}}

Ipotesi:

\begin{enumerate}
\def\labelenumi{\arabic{enumi}.}
\tightlist
\item
  completamento cicli L1-L2-L3 su claim prioritari;
\item
  repliche indipendenti multi-dominio;
\item
  governance stabile con KPI entro soglia.
\end{enumerate}

Tesi:
OCT può transitare da baseline fondativa a standard scientifico interoperabile.

Proof sketch:

\begin{enumerate}
\def\labelenumi{\arabic{enumi}.}
\tightlist
\item
  la robustezza teorica viene testata su domini eterogenei;
\item
  la robustezza empirica viene verificata da team indipendenti;
\item
  la robustezza istituzionale viene garantita da governance e audit.
\end{enumerate}

Conseguenze:
definisce criterio non arbitrario di maturità del programma.

Stato: \texttt{in\_proof}.

\subsection{Teorema 17.4 (Principio di cumulatività ordinativa)}

Ipotesi:

\begin{enumerate}
\def\labelenumi{\arabic{enumi}.}
\tightlist
\item
  pubblicazione integrale di successi e fallimenti;
\item
  aggiornamento versionato del ledger;
\item
  coerenza tra stato claim e rilascio pubblico.
\end{enumerate}

Tesi:
La conoscenza OCT cresce in modo cumulativo senza perdita di tracciabilità.

Proof sketch:

\begin{enumerate}
\def\labelenumi{\arabic{enumi}.}
\tightlist
\item
  la documentazione completa impedisce memoria selettiva;
\item
  il ledger versionato conserva storia decisionale;
\item
  l'allineamento release-state riduce distorsione comunicativa.
\end{enumerate}

Conseguenze:
fornisce base epistemica per fiducia inter-laboratorio.

Stato: \texttt{in\_proof}.

\begin{center}\rule{0.5\linewidth}{0.5pt}\end{center}

\section{7. Limiti finali e non-claim}

Limiti finali espliciti:

\begin{enumerate}
\def\labelenumi{\arabic{enumi}.}
\tightlist
\item
  il manoscritto e fondativo, non conclusivo;
\item
  alcuni blocchi formali richiedono dimostrazioni theorem-level più estese;
\item
  la validazione sociale necessità serie storiche e benchmark comparativi più ampi.
\end{enumerate}

Rischi aperti:

\begin{enumerate}
\def\labelenumi{\arabic{enumi}.}
\tightlist
\item
  inflazione terminologica senza corrispettivo empirico;
\item
  frammentazione tra repository/piattaforme se la governance si allenta;
\item
  rallentamento operativo per eccesso di proceduralità.
\item
  isolamento dalla letteratura adiacente per eccesso di riferimenti interni.
\end{enumerate}

Contromisure prioritarie:

\begin{enumerate}
\def\labelenumi{\arabic{enumi}.}
\tightlist
\item
  mantenere centralità dei quality gate e dei KPI;
\item
  pianificare cicli di semplificazione editoriale periodica;
\item
  promuovere confutazione esterna incentivata.
\item
  imporre sezione \texttt{related\ work} comparativa in ogni release \texttt{L2+}.
\end{enumerate}

Box C - Non-claim

Questo capitolo non implica:

\begin{enumerate}
\def\labelenumi{\arabic{enumi}.}
\tightlist
\item
  che OCT sia già standard globale;
\item
  che il programma sia immune da revisione profonda;
\item
  che la validazione futura confermera necessariamente tutti i claims correnti.
\end{enumerate}

\begin{center}\rule{0.5\linewidth}{0.5pt}\end{center}

\section{8. Claim Ledger (S0-S3)}

{\def\LTcaptype{none} % do not increment counter
\begin{longtable}[]{@{}lllll@{}}
\toprule\noalign{}
ID & Claim & Tipo & Evidenza & Stato \\
\midrule\noalign{}
\endhead
\bottomrule\noalign{}
\endlastfoot
C17.1 & L'opera 1-17 costituisce baseline fondativa pubblicabile & editoriale-scientifico & S1 & validated \\
C17.2 & La roadmap 2026-2030 e necessaria per transizione a standard & strategico-metodologico & S1 & validated \\
C17.3 & Il protocollo di adozione a livelli (\texttt{L0-L3}) e sufficiente come schema iniziale di rete & operativo & S1 & in\_review \\
C17.4 & La transizione a \texttt{v2.0} richiede evidenza inter-laboratorio, non solo interna & epistemico-istituzionale & S2 & in\_proof \\
C17.5 & La cumulatività ordinativa dipende da pubblicazione completa di successi/fallimenti & epistemico & S2 & in\_proof \\
C17.6 & La fruizione internazionale richiede doppio binario: didattica semplice + validazione rigorosa & dissemination & S1 & validated \\
\end{longtable}
}

\begin{center}\rule{0.5\linewidth}{0.5pt}\end{center}

\section{9. Chiusura}

OCT nasce in questo manoscritto come proposta forte ma disciplinata: forte nella pretesa di estendere la teoria delle categorie verso la composizione reale, disciplinata nel dichiarare limiti, stati e condizioni di verifica.

Il risultato principale dell'opera non e ``aver chiuso la teoria'', ma aver aperto una grammatica di ricerca in cui:

\begin{enumerate}
\def\labelenumi{\arabic{enumi}.}
\tightlist
\item
  la forma matematica dialoga con la prova empirica;
\item
  la validazione e tracciabile;
\item
  la conoscenza e cumulativa e contestabile.
\end{enumerate}

La prossima responsabilità non e scrivere slogan, ma costruire benchmark, repliche e comunità scientifica capace di migliorare anche contro l'autore.

Questa e la condizione reale perché OCT entri nella grammatica della scienza umana.

\begin{center}\rule{0.5\linewidth}{0.5pt}\end{center}

\section{Riferimenti}

Riferimenti di framework interno:

\begin{enumerate}
\def\labelenumi{\arabic{enumi}.}
\tightlist
\item
  Ghioni, F. (2026). \texttt{TE\_CORE\_v5.1.md}.
\item
  Ghioni, F. (2026). \texttt{Ordinative\_Set\_Theory\_OST\_A\_Concise\_Guide\_For\_AI\_v2\_1.md}.
\item
  \texttt{OCT\_FOUNDATIONAL\_BOOK\_CHAPTER\_01\_v1\_0.md} \ldots{} \texttt{OCT\_FOUNDATIONAL\_BOOK\_CHAPTER\_16\_v1\_0.md}.
\item
  \texttt{OCT\_VALIDATION\_PROTOCOL\_v0\_1.md}.
\item
  \texttt{DECISION\_MATRIX\_FINAL\_UNIFIED\_v0\_1.md}.
\item
  \texttt{OCT\_PUBLICATION\_PROGRESS\_TRACKER\_v0\_1.md}.
\end{enumerate}

Riferimenti primari:

\begin{enumerate}
\def\labelenumi{\arabic{enumi}.}
\tightlist
\item
  Popper, K. (1959). \emph{The Logic of Scientific Discovery}.
\item
  Lakatos, I. (1978). \emph{The Methodology of Scientific Research Programmes}.
\item
  Merton, R. K. (1973). \emph{The Sociology of Science}.
\item
  Letteratura contemporanea su reproducibility, preregistration e open scientific workflows.
\item
  Fong, B., Spivak, D. (2019). \emph{Seven Sketches in Compositionality}.
\item
  Baez, J., Erbele, J. (2015). Categories in control.
\item
  Coecke, B., Kissinger, A. (2017). \emph{Picturing Quantum Processes}.
\item
  Letteratura su open systems, resource theories e resilience quantitativa.
\end{enumerate}
